
\subsubsection{3.1 Data Collection and Coverage}

We use the Papers With Code archive (HuggingFace \texttt{pwc-archive/papers-with-abstracts}), filtering to papers appearing at NeurIPS, ICML, or ICLR during 2018--2024. This yields 12,600 papers with task annotations. Task labels serve as a proxy for benchmark datasets, as Papers With Code categorizes papers by the evaluation tasks they address.

For each paper, we record the set of tasks used for evaluation, venue, and year. Citation relationships are operationalized through task co-occurrence: papers evaluating on the same tasks and published in different years form proxy citation pairs, as papers on the same task naturally cite each other.

\subsubsection{3.2 Concentration Operationalization}

We measure benchmark concentration using the Herfindahl-Hirschman Index. For a given venue-year, let N be the number of papers and let n\_d be the count of papers using task d. The share for task d is s\_d = n\_d / N. HHI is:

$$\text{HHI} = \sum_d s_d^2$$

This metric equals the probability that two papers drawn uniformly at random from the venue-year evaluate on at least one common task. HHI ranges from 1/D (uniform distribution across D tasks) to 1 (all papers use the same task). Our observed HHI values range from 0.007 to 0.046.

We additionally compute normalized entropy as a diversity measure:

$$\text{Entropy} = -\sum_d p(d) \log p(d)$$

where p(d) is the proportion of papers using task d.

\subsubsection{3.3 Analytical Framework}

Our analysis traces the mechanism chain from concentration measurement to temporal dynamics:

\begin{enumerate}
\item \textbf{Validity}: Verify HHI can be computed across all venue-years and correlates with simpler metrics
\item \textbf{Propagation}: Test whether prior-year HHI predicts individual paper benchmark adoption
\item \textbf{Citation mechanism}: Test whether citing papers share benchmarks at higher rates than random pairs
\item \textbf{Temporal lock-in}: Test whether HHI in year t-1 reduces diversity in year t
\end{enumerate}

\subsubsection{3.4 Panel Regression Specification}

The temporal lock-in test (H-M5) uses the specification:

$$\text{Entropy}_{v,t} = \alpha + \beta \cdot \text{HHI}_{v,t-1} + \gamma_v + \delta_t + \epsilon_{v,t}$$

where $\gamma _v$ represents venue fixed effects and $\delta _t$ represents year fixed effects. A significant negative $\beta$ would indicate lock-in. We additionally conduct Granger causality tests.

